\documentclass[11pt,a4paper]{article}
\usepackage[utf8]{inputenc}
\usepackage[vietnamese]{babel}
\usepackage{amsmath,amssymb}
\usepackage{geometry}
\geometry{a4paper, total={180mm,260mm}, left=15mm, right=15mm, top=18mm, bottom=18mm}

\usepackage{tikz}
\usetikzlibrary{shapes.geometric, arrows.meta, positioning, calc}
\tikzset{
  arr/.style={-{Stealth[scale=0.8]}},
  every node/.append style={align=center}
}
\usepackage{tcolorbox}
\tcbuselibrary{skins,breakable}
\usepackage{fontawesome}
\usepackage{fancyhdr}
\usepackage{enumitem}

\pagestyle{fancy}
\fancyhf{}
\renewcommand{\headrulewidth}{0.5pt}
\renewcommand{\footrulewidth}{0.5pt}
\fancyhead[L]{\small\bfseries\color{violet!85!black}\faBook\ SỔ TAY KIẾN THỨC TOÁN 2}
\fancyhead[R]{\small\bfseries\color{violet!85!black}Lớp Toán Thầy Lĩnh -- loptoanthaylinh.vn}
\fancyfoot[C]{\small\color{gray}-- Sổ tay kiến thức: Bài 26. Đường gấp khúc. Độ dài đường gấp khúc --}

\setlength{\parindent}{0pt}
\setlength{\parskip}{0pt}

\begin{document}

% --- BANNER BÀI HỌC ---
\begin{tcolorbox}[
  enhanced,
  colback=violet!8!white,
  colframe=violet!70!black,
  arc=3mm, boxrule=1.2pt,
  left=3mm, right=3mm, top=2mm, bottom=2mm,
  halign=center
]
  {\LARGE\bfseries\color{violet!90!black}\faBookmark\ BÀI 26: ĐƯỜNG GẤP KHÚC. ĐỘ DÀI ĐƯỜNG GẤP KHÚC}\\[1mm]
  {\itshape\footnotesize\color{black!75} Nhận biết đường gấp khúc và quy tắc tính độ dài đường gấp khúc trên 1 trang A4}
\end{tcolorbox}

\vspace{-2mm}

\noindent
\begin{minipage}[t]{0.49\textwidth}\vspace{0pt}
\begin{tcolorbox}[
  enhanced,
  colback=violet!4!white,
  colframe=violet!75!black,
  arc=2mm, boxrule=0.8pt,
  left=2mm, right=2mm, top=1.8mm, bottom=1.8mm,
  title={\footnotesize\bfseries\color{white}\faKey\ 1. Khái niệm đường gấp khúc}
]
{\footnotesize
\begin{itemize}\setlength{\itemsep}{1pt}\setlength{\leftskip}{-2.5mm}
  \item \textbf{Đường gấp khúc} gồm nhiều đoạn thẳng nối tiếp nhau mà không cùng nằm trên một đường thẳng.
  \item Đọc tên đường gấp khúc lần lượt theo thứ tự các điểm từ đầu đến cuối: Đường gấp khúc $ABCD$.
  \item \textbf{Các đoạn thẳng tạo thành:}\\
  Đường gấp khúc $ABCD$ gồm $3$ đoạn thẳng: $AB, BC$ và $CD$.
  \item Hai đoạn thẳng liền nhau có chung một điểm (gọi là đỉnh).
\end{itemize}
}
\end{tcolorbox}
\end{minipage}
\hfill
\begin{minipage}[t]{0.49\textwidth}\vspace{0pt}
\begin{tcolorbox}[
  enhanced,
  colback=teal!4!white,
  colframe=teal!75!black,
  arc=2mm, boxrule=0.8pt,
  left=2mm, right=2mm, top=1.8mm, bottom=1.8mm,
  title={\footnotesize\bfseries\color{white}\faSliders\ 2. Tính độ dài đường gấp khúc}
]
{\footnotesize
\begin{itemize}\setlength{\itemsep}{1pt}\setlength{\leftskip}{-2.5mm}
  \item \textbf{Quy tắc vàng:}\\
  Độ dài đường gấp khúc bằng \textbf{tổng độ dài các đoạn thẳng} tạo nên đường gấp khúc đó.
  $$\mathbf{\text{Độ dài } ABCD = AB + BC + CD}$$
  \item \textbf{Điều kiện bắt buộc:}\\
  Độ dài các đoạn thẳng phải \textbf{cùng đơn vị đo} trước khi cộng.
  \item Nếu khác đơn vị đo, phải đổi về cùng một đơn vị đo rồi mới tính tổng.
\end{itemize}
}
\end{tcolorbox}
\end{minipage}

\vspace{-1.5mm}
\begin{tcolorbox}[
  enhanced,
  colback=orange!4!white,
  colframe=orange!80!black,
  arc=2mm, boxrule=0.8pt,
  left=2mm, right=2mm, top=1.5mm, bottom=1.5mm,
  title={\footnotesize\bfseries\color{white}\faSitemap\ 3. Mô hình trực quan đường gấp khúc ABCD}
]
\centering
\begin{tikzpicture}[x=1.5cm, y=0.8cm]
  \coordinate (A) at (0,0);
  \coordinate (B) at (1.8,1.2);
  \coordinate (C) at (4,0.2);
  \coordinate (D) at (5.8,1.5);

  \draw[very thick, violet!80!black] (A) -- (B) node[midway, above left, font=\scriptsize\bfseries] {$AB = 4\text{ cm}$};
  \draw[very thick, teal!80!black] (B) -- (C) node[midway, below, font=\scriptsize\bfseries] {$BC = 5\text{ cm}$};
  \draw[very thick, orange!85!black] (C) -- (D) node[midway, above right, font=\scriptsize\bfseries] {$CD = 6\text{ cm}$};

  \filldraw[violet!90!black] (A) circle (2.5pt) node[left=2pt, font=\scriptsize\bfseries] {$A$};
  \filldraw[black!80] (B) circle (2.5pt) node[above=2pt, font=\scriptsize\bfseries] {$B$};
  \filldraw[black!80] (C) circle (2.5pt) node[below=2pt, font=\scriptsize\bfseries] {$C$};
  \filldraw[orange!90!black] (D) circle (2.5pt) node[right=2pt, font=\scriptsize\bfseries] {$D$};

  \node[rectangle, rounded corners=2pt, draw=red!80!black, fill=red!10, font=\scriptsize\bfseries] at (3, -0.7) {
    Độ dài đường gấp khúc $ABCD = 4 + 5 + 6 = \mathbf{15\text{ cm}}$
  };
\end{tikzpicture}
\end{tcolorbox}

\vspace{-1.5mm}
\noindent
\begin{minipage}[t]{0.49\textwidth}\vspace{0pt}
\begin{tcolorbox}[
  enhanced,
  colback=green!4!white,
  colframe=green!60!black,
  arc=2mm, boxrule=0.8pt,
  left=2mm, right=2mm, top=1.8mm, bottom=1.8mm,
  title={\footnotesize\bfseries\color{white}\faCheckCircle\ 4. Ví dụ thực hành chuẩn mực}
]
{\footnotesize
\textbf{Bài toán mẫu:}\\
Tính độ dài đường gấp khúc $MNPQ$, biết độ dài các đoạn thẳng là: $MN = 12\text{ cm}$, $NP = 15\text{ cm}$, $PQ = 20\text{ cm}$.\\
\textit{Bài giải:}\\
Độ dài đường gấp khúc $MNPQ$ là:
$$12 + 15 + 20 = 47\text{ (cm)}$$
Đáp số: $47\text{ cm}$.
}
\end{tcolorbox}
\end{minipage}
\hfill
\begin{minipage}[t]{0.49\textwidth}\vspace{0pt}
\begin{tcolorbox}[
  enhanced,
  colback=red!4!white,
  colframe=red!75!black,
  arc=2mm, boxrule=0.8pt,
  left=2mm, right=2mm, top=1.8mm, bottom=1.8mm,
  title={\footnotesize\bfseries\color{white}\faExclamationTriangle\ 5. Mẹo tránh bẫy trắc nghiệm}
]
{\footnotesize
\begin{itemize}\setlength{\itemsep}{1pt}\setlength{\leftskip}{-2.5mm}
  \item \textbf{Bẫy khác đơn vị đo:}\\
  Nếu đề cho $AB = 1\text{ dm}$ và $BC = 5\text{ cm}$, không được lấy $1 + 5 = 6$. Phải đổi $1\text{ dm} = 10\text{ cm}$ rồi tính $10 + 5 = 15\text{ cm}$!
  \item \textbf{Đếm số đoạn thẳng:}\\
  Đường gấp khúc có $n$ điểm thì có $n - 1$ đoạn thẳng (VD: $ABCD$ có 4 điểm thì có 3 đoạn).
\end{itemize}
}
\end{tcolorbox}
\end{minipage}


\end{document}
